\chapter{总结与展望}

本章对全文工作进行总结，概括主要创新点，归纳研究局限，并从若干研究问题展望未来工作。

\section{全文工作总结}

本文围绕"面向异构 GPU 集群的多模态大语言模型并行训练技术"这一主题，针对实际环境中普遍存在的设备异构与 MLLM 的模块异构、分阶段冻结训练等特性，从\emph{模型划分与流水线调度}与\emph{重计算优化}两个层面开展了研究，提出了 MMoH 与 HR-MMoH 两套方法，并在异构集群上进行了实验验证。下文概括主要创新点。

\subsection{主要创新点}

本文的主要创新点可归纳为以下两方面。

\textbf{（1）MMoH：面向异构 GPU 集群的多模态大语言模型高效混合并行框架。}

针对异构集群上 MLLM 训练时存在的负载不均衡、资源利用率低以及冻结场景下冗余计算与通信等问题，本文提出了 MMoH 框架。其创新点包括：\emph{需求驱动的设备拓扑识别}，根据 MLLM 各模块的计算与显存需求趋势与异构设备的算力—显存供给进行匹配，筛选候选设备拓扑，使需求与供给在趋势上一致。\emph{复合粒度模型抽象}，针对编码器、投影层与 LLM 基座的结构差异，采用多级粒度将模型重组为模型块列表（编码器按子模块合并、LLM 按注意力/FFN 子层拆分），在划分复杂度与负载均衡之间取得折中。\emph{冻结感知的模型划分}，将块到阶段的分配形式化为整数规划问题，在满足显存与主阶段约束下最小化迭代时间，且目标函数显式依赖各阶段的冻结状态，从而随训练阶段动态调整划分策略。\emph{提前终止调度器}，在 1F1B 等流水线调度基础上，识别从首端开始的连续冻结阶段并对其省略反向计算与流水级间通信，在不影响模型收敛性的前提下提升吞吐。在两种异构集群（Cluster-S、Cluster-L）与两种规模 MLLM（MLLM-3B、MLLM-12B）上的实验表明，MMoH 在 no-freeze、freeze-encoder、freeze-llm、freeze-both 四种场景下相对 Megatron-LM、Whale、Espresso 均可取得最高吞吐，最高可达约 1.77$\times$ 加速，且提前终止调度器在冻结场景下可带来约 8\%--19\% 的额外增益，并具有良好的可移植性。

\textbf{（2）HR-MMoH：面向异构设备与异构模型的重计算策略优化。}

针对现有重计算策略在同构假设下采用全局统一配置、难以在异构集群上同时兼顾显存安全与计算效率的问题，本文在 MMoH 的划分与调度基础上，提出并实现了 HR-MMoH 的异构重计算方法。HR-MMoH 使各流水级根据自身显存与算力条件采用差异化的重计算粒度（如 full/selective）与子模块选择，在降低 OOM 风险的同时减少多余重算、平衡流水级负载，与 MMoH 形成 "划分—调度—重计算" 的联合优化。第四章在更大批次大小与更长序列长度的实验下证明，HR-MMoH 在 no-freeze、freeze-encoder、freeze-llm、freeze-both 四种冻结场景中均取得最高吞吐，相对 Espresso 的迭代时间缩短约 6.1\%--23.1\%，在 Whale 可训练的三种冻结场景中缩短约 11.0\%--40.6\%。这说明按流水级差异化配置能够系统性优于全局统一策略，并验证了与 MMoH 协同优化的有效性。

综上，本文在异构 GPU 集群上 MLLMs 并行训练这一方向上，从\emph{划分与调度}与\emph{显存优化}两条线给出了系统性的方法设计与实验支撑，为实际部署中的混合集群与分阶段训练场景提供了可参考的技术路径。

\section{局限性及未来工作展望}

结合前文全文总结，本节先归纳研究局限，再围绕后续研究的核心问题与思路作简要展望。

\subsection{研究局限性}

本研究仍存在以下局限，亦是后续工作的重要出发点。

规划器中的整数规划求解在模型层数与设备数较大时可能面临组合爆炸，当前采用启发式或现有求解器在有限时间内求近似解，最优性保证与可扩展性仍有提升空间；实验仅在两种规模的 MLLM（3B、12B）与两种异构集群配置上进行，对更大规模模型（如 70B+）、更多代际与型号混合的集群以及更长序列、更多模态的 MLLMs 的泛化能力有待进一步验证。

尽管第四章已完成端到端实验并验证了按流水级异构重计算的收益，当前策略搜索仍以启发式与剖析驱动为主，候选配置空间、搜索顺序与回退机制依赖经验设计，尚未给出更强的最优性保证与复杂度分析；实验主要在 Cluster-L 与 MLLM-12B 上开展，对更大规模模型、更复杂模型结构（如 MoE 模型）及更长上下文设定下的鲁棒性仍需补充；此外，HR-MMoH 与 MMoH 的联合优化目前仍采用"先划分、后配置重计算"的分层流程，尚未形成统一目标下的联合求解框架。

当前工作主要针对视觉—语言类 MLLM 与 NVIDIA GPU 生态，对音频、视频等多模态输入、以及 AMD、国产 GPU 等异构硬件与软件栈的适配与评估尚未涉及；与专家并行、上下文并行等已有技术的协同优化仍有扩展空间。

\subsection{未来研究工作}

基于上述局限，后续研究可沿以下两条线索推进。

\textbf{（1）划分与重计算的联合求解。}
本文采用"先由 MMoH 求解划分与设备拓扑，再由 HR-MMoH 为各流水级搜索重计算配置"的分层流程。该流程存在循环依赖：划分决定各流水级的层分配与显存余量，进而影响可选的重计算配置与各流水级耗时；而重计算配置改变后各流水级的实际耗时又反过来改变最优划分。当前分层求解无法捕捉这种耦合，可能遗漏"不同划分 + 不同重计算组合"下迭代时间更低的方案。一种可行思路是将重计算配置作为决策变量引入第三章的整数线性规划，把重计算带来的额外前向开销纳入迭代时间目标，使划分与重计算在同一目标下联合求解；或采用交替优化——固定重计算求划分，再固定划分更新重计算——并分析其收敛性。此方向的核心在于量化分层流程的次优性来源及其量级，为判断是否值得引入联合求解提供依据。

\textbf{（2）提前终止调度向新型流水线机制的扩展。}
第三章的提前终止调度器与迭代时间的推导均建立在 1F1B 调度之上。近年提出的 ZeroBubble\cite{zerobubble2024} 将反向拆分为"对输入的梯度"与"对参数的梯度"两个独立阶段，通过 1F1B1W 调度将后者与下一 micro-batch 的前向重叠以压缩气泡；DualPipe\cite{dualpipe2025} 则采用双向 micro-batch 流动实现前反向交织。
这两类调度改变了流水级间的依赖结构：冻结流水级不需要计算对参数的梯度，但对输入的梯度能否被提前终止剪枝，取决于后续流水级在新调度下的依赖关系是否与 1F1B 一致。
在上述新调度框架下重新推导冻结前缀的可剪枝条件与迭代时间收益，并验证其与 HR-MMoH 按流水级差异化重算的兼容性可作为下一步研究工作。
若成功推广，可在进一步压缩气泡的同时保留冻结前缀剪枝的收益，使 MMoH 框架适用于更广泛的工程实践。

综上所述，本文在面向异构 GPU 集群的 MLLMs 并行训练策略方面取得了阶段性成果；后续宜在划分与重计算的联合求解、以及提前终止向新型流水线调度的扩展两个方向上继续推进，为异构环境下 MLLMs 训练的进一步优化提供方法与理论支撑。
